% Generated draft: requires mathematical and source review.
\documentclass[11pt]{article}
\usepackage[margin=25mm]{geometry}
\usepackage{fontspec}
\setmainfont{Latin Modern Roman}
\usepackage{amsmath,amssymb,mathtools,unicode-math}
\setmathfont{Latin Modern Math}
\usepackage{enumitem,needspace,xurl,hyperref}
\setlength{\parindent}{0pt}
\setlength{\parskip}{5pt}
\setlength{\emergencystretch}{3em}
\newcommand{\N}{\mathbb{N}}
\newcommand{\Z}{\mathbb{Z}}
\newcommand{\Q}{\mathbb{Q}}
\newcommand{\R}{\mathbb{R}}
\newcommand{\C}{\mathbb{C}}
\newcommand{\F}{\mathbb{F}}
\newcommand{\A}{\mathbb{A}}
\newcommand{\PP}{\mathbb{P}}
\newcommand{\E}{\mathbb{E}}
\newcommand{\eps}{\varepsilon}
\newcommand{\dd}{\,\mathrm{d}}
\begin{document}
\section*{UnsolvedMath ID: 30000865}
\textbf{Source problem code:} catalogue entry
\subsection{Definitions and mathematical statement}
Let $p(n)$ be the number of partitions of the nonnegative integer $n$, with $p(0)=1$. For $z\in\mathbb H=\{z\in\C:\operatorname{Im}z>0\}$ put $q=e^{2\pi iz}$, and interpret $q^{1/24}=e^{2\pi iz/24}$. The supplied identity is
\[
P(q)=\sum_{n\geq0}p(n)q^n
=1+\sum_{n=1}^{\infty}\frac{q^{n^2}}{\prod_{j=1}^{n}(1-q^j)^2}
=\frac{q^{1/24}}{\eta(z)},
\qquad
\eta(z)=q^{1/24}\prod_{j\geq1}(1-q^j).
\]
The source question is to prove directly from this hypergeometric expansion that $P(q)$ is modular; the original wording says that it is a modular form.

There is a normalization issue in that wording. The usual modular object associated with the displayed identity is
\[
H(z)=q^{-1/24}P(q)=\eta(z)^{-1},
\]
rather than the unnormalized $P(q)$. A possible precise normalized target would be to derive directly from the hypergeometric series the transformation laws
\[
H(z+1)=e^{-\pi i/12}H(z),\qquad
H(-1/z)=(-iz)^{-1/2}H(z),
\]
with the holomorphic square root positive when $z$ is positive imaginary, together with the appropriate cusp behavior. In the conventional terminology this normalized object has weight $-1/2$, the inverse eta multiplier, and is weakly holomorphic, allowing a pole at a cusp. This is an explicitly qualified interpretation of the question, not a silently corrected replacement of its claim about $P$. The supplied record does not specify weight, group, multiplier, cusp convention, or whether normalization by $q^{-1/24}$ is intended. The allowed meaning of a proof ``directly'' from the hypergeometric expansion also needs specification; merely citing the displayed eta identity and known eta transformations would not by itself address that requested method.

The registry labels the statement corrected and verified, saying that the full identity and direct-modularity question were restored from the official report. Its literature triage, checked 2026-08-21, labels this particular direct-proof question open and says no later resolving direct proof was identified; these are supplied assessments.
\subsection{Short English statement}
Derive the modular behavior of the partition series directly from its displayed hypergeometric expansion. The source calls $P$ modular, but the usual modular form is its normalized version $q^{-1/24}P=\eta^{-1}$; the intended normalization and meaning of a direct proof require review.
\subsection{Sources}
\begin{itemize}
\item Local UnsolvedMath registry, permanent ID 30000865, source code catalogue entry; supplied complete record, including its corrected-verified statement note and dated literature triage.
\item Modulformen (2008), Oberwolfach Reports, DOI 10.4171/owr/2007/51, pp. 3013--3072: \url{https://doi.org/10.4171/owr/2007/51}. Proposed-year metadata gives 2007.
\item Supplied report PDF-text reference: \url{https://oa.tib.eu/renate/bitstreams/89e1927e-e09f-44be-b490-7a32b587d287/download}.
\end{itemize}
\end{document}
